\section{Первая часть курса.}
\subsection{Домашние задачи на первую неделю}
\begin{problem}
    В выпуклом $n$-угольнике никакие три диагонали не пересекаются в одной точке внутри этого $n$-угольника. Случайным образом выбираются две различные его диагонали. Найдите вероятность, что они пересекаются внутри $n$-угольника. Найдите предел этой вероятности при $n \rightarrow +\infty$.
\end{problem}
\begin{solution}
    Будем считать, что $n \geq 4$, поскольку иначе в многоугольнике в принципе нет диагоналей. \\
    Рассмотрим две какие-то диагонали, пересекающиеся внутри многоугольника. Заметим, что они однозначно определяют способ выбрать четыре вершины многоугольника. \\
    Теперь рассмотрим произвольные 4 различные точки многоугольника. Из выпуклости следует, что существуют единственные две пересекающиеся внутри многоугольника диагонали, вершинами которых являются эти четыре точки. Таким образом, количество способов выбрать две пересекающиеся диагонали в выпуклом $n$-угольнике равно $\binom{n}{4} = \frac{n(n -1)(n - 2)(n - 3)}{4!}$. Поскольку всего в многоугольнике $\frac{n(n - 3)}{2}$ диагоналей, то всего способов выбрать две диагонали -- $\frac{n(n - 3)}{2}\bigg(\frac{n(n - 3)}{2} - 1\bigg)\cdot \frac{1}{2}$. Таким образом, итоговая вероятность $\mathbb{P}(n)$ выражается как \[
                                                                                                                                                                                                                                                                                                                                                                                                                                                                                                                                                                                                                                                \mathbb{P}(n) = \dfrac{\dfrac{n(n - 1)(n - 2)(n - 3)}{24}}{\dfrac{n(n - 3)(n(n - 3) - 2)}{8}} = \dfrac{1}{3}\cdot \dfrac{n^2 - 3n + 2}{n^2 - 3n - 2}
    \]
    В пределе при $n \rightarrow +\infty$: \[
                                               \mathbb{P}(n) \rightarrow \dfrac{1}{3}, n \rightarrow +\infty
    \]
\end{solution}
\begin{problem}
    При двух бросках игральной кости выпало $X_1$ и $X_2$, соответственно. Вычислите сумму средних значений величин $\max\{X_1, X_2\}$ и $\min\{X_1, X_2\}$.
\end{problem}
\begin{solution}
    \begin{multline*}
        \mathbb{E}\max\{X_1, X_2\} + \mathbb{E}\min\{X_1, X_2\} = \int_\Omega \max\{X_1, X_2\}d\mathbb{P} + \int_\Omega \min\{X_1, X_2\}d\mathbb{P} =  \\ = \int_\Omega \max\{X_1, X_2\} + \min\{X_1, X_2\}d\mathbb{P} = \int_\Omega X_1 + X_2d\mathbb{P} = \\ = \int_\Omega X_1d\mathbb{P} + \int_\Omega X_21d\mathbb{P} = 2 \cdot \sum\limits_{i = 1}^6 \dfrac{i}{6} = 7
    \end{multline*}
\end{solution}
\begin{problem}
    Дан клетчатый прямоугольник высоты 4 и ширины $N$. Вася красит какой-то горизонтальный прямоугольник $1 \times 3$ клетки, а Петя красит какой-то вертикальный прямоугольник $1 \times 3$ клетки. Найдите вероятность того, что хотя бы одна клетка будет покрашена дважды.
\end{problem}
\begin{solution}
    Будем считать, что $N \geq 3$. \\
    У Васи есть всего $4 \cdot (N - 2)$ возможных выборов прямоугольников. У Пети есть $(4 - 2)\cdot N = 2N$  возможных выборов прямоугольников. Тогда, пусть Вася выбрал один из $2(N - 2)$ прямоугольников в крайних строках. Для каждого из них у Пети есть ровно три способа выбрать свой прямоугольник так, чтобы он пересекался с Васиным. Если же Вася выберет один из $2(N - 2)$ средних прямоугольников, то каждый из них пересекается с 6 горизонтальных прямоугольников. Таким образом, всего пар (горизонтальных $1 \times 3$, вертикальный $1 \times$), которые пересекаются по какой-нибудь клетке ровно $18(N - 2)$. Значит, вероятность того, что хотя бы одна клетка будет покрашена дважды $\mathbb{P}(N)$ равна \[
                                                                                                                                                                                                                                                                                                                                                                                                                                                                                                                                                                                                                                                                                                                                       \mathbb{P}(N) = \dfrac{18(N - 2)}{4(N - 2)\cdot 2N} = \dfrac{9}{4N}
    \]
\end{solution}
\begin{problem}[ВТФ]
    Две урны содержат одинаковое количество шаров. Шары окрашены в белый и черный цвета. Из каждой урны вынимают по $n$ шаров с возвращением, где $n \geq 3$. Найдите $n$ и «состав» каждой урны, если вероятность того, что все шары, взятые из первой урны, белые, равна вероятности того, что все шары, взятые из второй урны, либо белые, либо черные.
\end{problem}
\begin{solution}
    Пусть в первой урне $x$ белых шаров, во второй -- $y$. \\
    Вероятность того, что все шары из первой урны -- белые, равна $\bigg(\frac{x}{n}\bigg)^n$. \\
    Вероятность того, что все шары из второй урны -- белые, или все шары из второй урны -- чёрные равна $\bigg(\frac{y}{n}\bigg)^n + \bigg(\frac{n - y}{n}\bigg)^n$. \\
    По условию они равны. Тогда \[
                                    x^n = y^n + (n - y)^n, n \geq 3
    \]
    Согласно Великой Теореме Ферма, это уравнение не имеет нетривиальных решений при $n \geq 3, x, y \in \N$. Значит, или $y = 0$, или $n - y = 0$, т.е. $x = n, y = 0$ или $x = n, y = n$.
\end{solution}
\begin{problem}
    Симметричную монетку бросают неограниченное число раз. Какая из последовательностей встретится раньше с большей вероятностью: $POP$ или $PPO$?
\end{problem}
\begin{solution}
    Пусть в некоторый момент встретилась в первый раз буква $P$ (до этого выпадала только $O$). Пусть $POP$ встречается раньше $PPO$ с вероятностью $\alpha$. Заметим, что если в некоторый момент не встретилось $PPO$ и последние два -- $OO$, то для строки, начинающейся с $OO$ вероятность встретить $POP$ раньше $PPO$ тоже будет равна $\alpha$. Тогда, рассмотрим первую встреченную $P$ и рассмотрим несколько возможных исходов:
    \begin{equation*}
        P = \begin{cases}
                OO, & \alpha/4 \\
                OP, & 1/4 \\
                PO,  & 0 \\
                PP = 1/4\cdot\begin{cases}
                                 OO, & 0 \\
                                 OP, & 0 \\
                                 PO, & 0 \\
                                 PP, & \ldots
                \end{cases}
        \end{cases}
    \end{equation*}
    То есть \[
                \alpha = \dfrac{1}{4}(\alpha + 1) + \dfrac{1}{4}\bigg(0 + \dfrac{1}{4}\bigg(0 + \ldots\bigg)\bigg)
    \]
    Значит \[
               \alpha = \dfrac{1}{4}(\alpha + 1) \Rightarrow \alpha = \dfrac{1}{3}
    \]
    Поскольку вероятность того, что не встретится ни $POP$, ни $PPO$ равна 0, то вероятность того, что $PPO$ встретится раньше равна $\frac{2}{3}$.
\end{solution}